\chapter{Introduction and Planning}
\label{chapter:introduccion_and_planning}

\section{Introduction}

Presentación general del trabajo (Simulación de riadas).

\section{Motivation and Justification}

Por qué es importante este trabajo.

\section{Objectives}

Objetivos generales y específicos.

\section{Project Planning and Management}
\label{sec:project_planning_and_management}

The successful execution of the Final Master's Project (FMP), which focuses on developing a scientific C++ program for flood simulation, requires robust planning and diligent management. This course is weighted at 18 ECTS, equating to an estimated dedicated effort of approximately 450 hours. The timeline presented below serves as a guiding roadmap, designed to track progress against key project milestones and deliverables.

\subsection{Work Methodology}
\label{subsec:methodology}

The project methodology is structured around the four mandatory phases defined for the FMP, integrated with a phased software development approach. This hybrid methodology ensures that the core theoretical and design elements (Phases 1 \& 2) are completed before critical implementation and validation (Phase 3) and final documentation (Phase 4). The iterative nature of the plan allows for necessary adjustments and includes dedicated time for thorough testing and writing throughout the academic year.

\subsection{Project Schedule: Gantt Chart}
\label{subsec:gantt}

The following Gantt chart (Figure \ref{fig:ganttchart}) illustrates the proposed 24-week timeline for the project, distributing the estimated 450 working hours across the defined phases. Time buffers have been implicitly integrated into the task durations to account for unforeseen issues, necessary code refactoring, and academic holidays.

% ----------------------------------------------------------------
% START OF GANTT CHART CODE
% ----------------------------------------------------------------

\begin{figure}[h!]
    \resizebox{\linewidth}{!}{
        \centering
        \begin{ganttchart}[
            vgrid,
            bar height=0.6,
            group height=0.3,
            title label font=\bfseries\normalsize
        ]{1}{24}
            \gantttitle{Project Timeline (24 Weeks)}{24} \\
            \gantttitlelist{1,...,24}{1} \\
            % PHASE 1
            \ganttgroup[group label node/.append style={fill=blue!20}]{Phase 1: Planning (4 Weeks)}{1}{4} \\
            \ganttbar[bar/.append style={fill=blue!50}]{P1.1 State-of-the-Art Review}{1}{3} \\
            \ganttbar[bar/.append style={fill=blue!50}]{P1.2 Objectives and Scope Definition}{3}{4} \\
            \ganttbar[bar/.append style={fill=blue!50}]{P1.3 Development Environment Setup}{1}{1} \\ % Etiqueta acortada
            % PHASE 2
            \ganttgroup[group label node/.append style={fill=green!20}]{Phase 2: Problem \& Proposal (6 Weeks)}{5}{10} \\
            \ganttbar[bar/.append style={fill=green!50}]{P2.1 Cellular Automata (CA) Design}{5}{6} \\
            \ganttbar[bar/.append style={fill=green!50}]{P2.2 GIS Data Reading Implementation}{5}{7} \\ % Etiqueta acortada
            \ganttbar[bar/.append style={fill=green!50}]{P2.3 Water Propagation Rules (CA)}{7}{10} \\
            % PHASE 3
            \ganttgroup[group label node/.append style={fill=yellow!20}]{Phase 3: Partial Deliverable (6 Weeks)}{11}{16} \\
            \ganttbar[bar/.append style={fill=yellow!70}]{P3.1 Unit Testing of the CA Engine}{11}{12} \\
            \ganttbar[bar/.append style={fill=yellow!70}]{P3.2 Integration and Data Validation}{13}{15} \\
            \ganttbar[bar/.append style={fill=yellow!70}]{P3.3 Progress Document Generation. (Part 1)}{14}{16} \\ % Etiqueta acortada
            % PHASE 4
            \ganttgroup[group label node/.append style={fill=red!20}]{Phase 4: Final Deliverable (6 Weeks)}{17}{22} \\
            \ganttbar[bar/.append style={fill=red!50}]{P4.1 Results Analysis and Visualization}{17}{19} \\
            \ganttbar[bar/.append style={fill=red!50}]{P4.2 Final Thesis Writing and Revision}{18}{21} \\
            \ganttbar[bar/.append style={fill=red!50}]{P4.3 Oral Presentation Preparation}{20}{22} \\
            % PHASE 5
            \ganttgroup[group label node/.append style={fill=gray!20}]{Phase 5: Presentation and Defense (2 Weeks)}{23}{24} \\
            \ganttbar[bar/.append style={fill=gray!50}]{P5.1 Final Supervisor Review}{23}{23} \\
            \ganttbar[bar/.append style={fill=gray!50}]{P5.2 Virtual Defense}{24}{24}
        \end{ganttchart}
    }
    \caption{Project Timeline (24 Weeks)}
    \label{fig:ganttchart}
\end{figure}

% ----------------------------------------------------------------
% END OF GANTT CHART CODE
% ----------------------------------------------------------------

\clearpage

\subsection{Details of Activities}
\label{subsec:activity_details}

The following table provides a detailed breakdown of the tasks defined in the timeline, specifying the scope of work for each activity.

% ----------------------------------------------------------------
% START OF LONGTABLE CODE
% ----------------------------------------------------------------

\begin{longtable}{p{0.1\linewidth} p{0.2\linewidth} p{0.15\linewidth} p{0.5\linewidth}}
    \caption{Detailed Breakdown of Final Master's Project Activities.} \\
    \toprule
    \textbf{ID} & \textbf{Activity} & \textbf{FMP Phase} & \textbf{Description (2-5 Lines)} \\
    \midrule
    \endhead

    % PHASE 1
    \multicolumn{4}{l}{\textbf{PHASE 1: PLANNING}} \\
    \midrule
    P1.1 & State-of-the-Art Review & Phase 1 & Review existing scientific literature on flood simulation, focusing on Cellular Automata (CA) and C++ flow models. Key CA models suitable for water physics will be identified and the core references for the thesis will be collected. \\
    P1.2 & Objectives and Scope Definition & Phase 1 & Clearly specify the flood scenarios the program will simulate (e.g., 2D vs. 3D, terrain type). Success criteria will be established, and the exact GIS data format required for input topography will be defined. \\
    P1.3 & Development Environment Setup & Phase 1 & Install and configure the C++ compiler (including necessary GIS libraries like GDAL), the version control system (Git), and development tools. A basic LaTeX project structure will be created to initiate the project documentation. \\
    \midrule

    % PHASE 2
    \multicolumn{4}{l}{\textbf{PHASE 2: PROBLEM DESCRIPTION AND PROPOSAL}} \\
    \midrule
    P2.1 & Cellular Automata (CA) Design & Phase 2 & Define the transition rules and neighborhood structure for the CA (e.g., Moore or von Neumann). The C++ data structure will be designed to efficiently represent terrain cells and water depth. \\
    P2.2 & GIS Data Reading Implementation & Phase 2 & Develop the C++ module for correctly reading and parsing the input GIS data (Digital Elevation Model or DEM). This module must initialize the initial state of the Cellular Automata based on the topography and the flood starting point. \\
    P2.3 & Water Propagation Rules (CA) & Phase 2 & Code the core of the simulator, which includes the equations and logic for water transfer between cells in each iteration. This represents the underlying flow model that governs the flood simulation. \\
    \midrule

    % PHASE 3
    \multicolumn{4}{l}{\textbf{PHASE 3: PARTIAL DELIVERABLE}} \\
    \midrule
    P3.1 & Unit Testing of the CA Engine & Phase 3 & Conduct exhaustive testing of critical code functions, ensuring that GIS data reading, terrain initialization, and, mainly, the CA water propagation rules behave as expected. Focus will be on identifying precision and performance errors. \\
    P3.2 & Integration and Data Validation & Phase 3 & Run the complete simulator with a representative set of GIS data. Simulation results will be compared against known flood data or other model outputs to validate the program's behavior. \\
    P3.3 & Progress Document Generation (Part 1) & Phase 3 & Compile and draft the initial chapters of the FMP document in LaTeX. This includes the introduction, literature review, and a detailed description of the development methodology and the Cellular Automata design. \\
    \midrule

    % PHASE 4
    \multicolumn{4}{l}{\textbf{PHASE 4: FINAL DELIVERABLE}} \\
    \midrule
    P4.1 & Results Analysis and Visualization & Phase 4 & Process the simulation output data (depth maps, flooding time, etc.). External tools or C++ libraries will be used to generate clear graphs and visualizations that allow the interpretation of the flood impact. \\
    P4.2 & Final Thesis Writing and Revision & Phase 4 & Complete the Results and Conclusions chapters of the LaTeX document. A full review of style, grammar, and formatting will be performed to ensure compliance with academic standards. \\
    P4.3 & Oral Presentation Preparation & Phase 4 & Design and create the presentation slides. The speech will be rehearsed, emphasizing the justification of the CA model, the C++ development process, and the key results obtained with the GIS data. \\
    \midrule
    
    % PHASE 5
    \multicolumn{4}{l}{\textbf{PHASE 5: PRESENTATION AND VIRTUAL DEFENSE}} \\
    \midrule
    P5.1 & Final Supervisor Review & Phase 5 & Final session with the FMP supervisor to address last-minute corrections to the document, code, and presentation. Minor details will be adjusted before the final submission and defense. \\
    P5.2 & Virtual Defense & Phase 5 & Formal presentation of the FMP to the evaluation committee. Technical decisions, methodology, and project conclusions will be defended. \\
    \bottomrule
\end{longtable}

% ----------------------------------------------------------------
% END OF LONGTABLE CODE
% ----------------------------------------------------------------


\section{Document Structure}

Breve resumen de los contenidos de los siguientes capítulos.


